\section{Compression Trees and Local Laws}\label{sec:min-framework}

Let $\Strings$ denote the set of finite token sequences with concatenation
written $u \concat v$. The task target is an oracle
$\fstar_{\mathrm{econ}} : \Strings \to \Y$ with discrepancy metric
$\metric$. A summarizer is a (possibly stochastic) map
$g : \Strings \rightsquigarrow \Strings$---a prompted language model in the
Manifesto application, a learned neural operator in the Markov anchor, a
serialized hand-designed register array in the HLL anchor. A readout
$f : \Strings \rightsquigarrow \Y$ scores the root summary: in closed-form
examples $f$ may be taken to be $\fstar_{\mathrm{econ}}$ itself, and in the
Manifesto results $f$ is a rubric scorer that approximates
$r_{\mathrm{econ}}$---the rubric scorer named in
\S\ref{sec:min-manifesto-state}. Every call to $g$ is \emph{local}: its
input is either a leaf block or the concatenation of two child summaries.
That locality is what makes each check small enough to audit, and what
keeps the guarantees conditional on claims one node at a time rather than
on the whole document at once.

A document $x$ is partitioned into contiguous leaves $(b_1, \ldots, b_k)$
attached to the leaves of a full binary tree $T$. Each leaf stores
$s_i = g(b_i)$, and each internal node $u$ stores
$s_u = g(s_{u_L} \concat s_{u_R})$. Only these stored summaries move
through the pipeline at deployment; the raw span underneath a node is
retained conceptually, so that a later audit can compare the stored state
to the evidence it is supposed to preserve.
Figure~\ref{fig:min-plain-tree} shows the object. Observation itself can
happen at two tiers: the root-observed view already introduced in
\S\ref{sec:min-manifesto-state}, where the tree root is compared directly
to a full-document expert target; and a node-level local-law audit that
observes or judges selected leaves, stored summaries, and merge nodes
inside the realized tree. The theorems below do not mix the two. They
make preservation claims about the tree that either view can cash in on
its own terms.

\begin{figure}[t]
    \centering
    \includegraphics[width=0.92\linewidth]{01_base_plain.pdf}
    \caption{\textbf{A compression tree reduces a document to one root
    summary by recursively merging adjacent child summaries.} Leaves
    summarize raw spans; internal nodes merge adjacent child summaries;
    only summaries are retained on each node. In the Manifesto running
    example, every stored string is treated as a surrogate for the
    economic-policy state of the span beneath it.}
    \label{fig:min-plain-tree}
\end{figure}

Two strings are \emph{oracle-equivalent}, written
$u \sim_{\mathrm{econ}} v$, when
$\metric(\fstar_{\mathrm{econ}}(u), \fstar_{\mathrm{econ}}(v)) = 0$. All
local laws are stated modulo this relation, which is worth underlining:
``preserving'' a span does not mean producing a string that resembles it
under some generic semantic metric, but producing a string that the
economic oracle cannot tell apart from it. Two paragraphs with very
different surface form can be $\sim_{\mathrm{econ}}$, and two paragraphs
with very similar surface form can fail to be equivalent if one tucks in
a fiscal qualifier the other drops. For the induction to carry a local
equivalence up the tree, one more property is needed: equivalence has to
survive insertion into the same surrounding context.

\begin{assumption}[Context-compatible oracle]\label{ass:context}
For every context $w \in \Strings$, $u \sim_{\mathrm{econ}} v$ implies
both $w \concat u \sim_{\mathrm{econ}} w \concat v$ and
$u \concat w \sim_{\mathrm{econ}} v \concat w$.
\end{assumption}

This is a task assumption, not a universal property of text. It is
plausible for rubrics that read evidence locally---the economic dimension,
a clinical adverse-event count, a citation-extraction pipeline. It is less
plausible, and should be treated as a source of auditable distortion, for
discourse-sensitive targets like irony or rhetorical framing, where
substituting an oracle-equivalent span can change how neighboring
paragraphs are read.

The three local laws are the auditable conditions under which a summary at
any node actually preserves the oracle-visible state of the span below it.

\paragraph{Condition 1: Sufficiency (C1).}
\begin{equation}\tag{C1}\label{eq:leaf_sufficiency}
    \forall\text{ realized leaf } b : \qquad
    g(b) \sim_{\mathrm{econ}} b.
\end{equation}
A leaf summary must keep the evidence the economic oracle would read in
its span. C1 fails when a paragraph summary drops a decisive policy
commitment, a fiscal qualifier, or an entity whose presence the rubric
would score.

\paragraph{Condition 2: Idempotence (C2).}
\begin{equation}\tag{C2}\label{eq:leaf_idempotence}
    \forall\text{ stored summary } s \in \operatorname{range}(g) : \qquad
    g(s) \sim_{\mathrm{econ}} s.
\end{equation}
A second summarization pass over an already-stored summary must not shift
the economic readout. C2 fails when a clean-up pass smooths away a
decisive qualifier---an already-reduced summary can be reduced further
into text so generic that it no longer places the platform on the rubric.
When $g$ is stochastic, take $\operatorname{range}(g)$ to be the
support of its output PMF and read $\sim_{\mathrm{econ}}$ in
expectation over the draws; Appendix~\ref{app:min-proofs} carries the
PMF formulation used in the proofs.

\paragraph{Condition 3: Merge Consistency (C3).}
\begin{equation}\tag{C3}\label{eq:leaf_compatibility}
    \forall\, u, v \in \Strings : \qquad
    u \concat v
    \sim_{\mathrm{econ}} g(u \concat v)
    \sim_{\mathrm{econ}} g\!\left(g(u) \concat g(v)\right).
\end{equation}
The first equivalence says a joint summary of two spans preserves the
economic reading of their union. The second says that summarizing two
children separately and merging the summaries preserves the same
cross-span interaction as summarizing the raw concatenation at once---a
platform that is environmentally interventionist in one section but
fiscally restrained in another has its economic reading at the boundary
between those sections, and that reading must survive the merge. C3 is
the law the merge has to satisfy for the tree induction to run.

A \emph{global} law is the universal statement above, quantified over all
$b$, $s$, $u$, $v$. A \emph{realized} law is the corresponding check
instantiated on the specific strings produced in one execution of the
tree. The audit never proves the global form; it estimates whether the
deployed run is realized-valid often enough to support the guarantee the
method is claiming.

\paragraph{Assumption bundles.}
Four named bundles are cited in the theorems of
Section~\ref{sec:min-guarantees}.
\emph{Preservation} = C1 + C3 + Assumption~\ref{ass:context}; it delivers
one-pass root preservation. \emph{Recompression} = Preservation + C2; it
extends preservation through any number of additional re-summarization
rounds. \emph{Optimization} = Recompression + a factorization /
measurability premise on the downstream objective; it makes training on
root summaries and training on full documents share a population target.
\emph{Certificate} = Optimization + transport, calibration, and sampling
assumptions; it converts sampled audit statistics into a finite-sample
objective-gap certificate. The first two are structural assumptions about
the compression object; the last two add downstream modelling ingredients.
Appendix~\ref{app:min-notation} gives the explicit table with the exact
theorem each bundle backs, and Appendix~\ref{app:min-proofs} records the
C1 / C3 / C2 $\leftrightarrow$ L1 / L3 / L2 crosswalk to the Lean
formalization.
